% =====================================================================
%  IEL04 - Circuitos Eléctricos I
%  Semana 8: Circuitos RL en serie y en paralelo: reactancia inductiva, impedancia y leyes de Kirchhoff con la bobina real
%  Institución Universitaria de Barranquilla (IUB)
%  Generada por src/presentaciones.py (diseño IUB 5). Compilador: pdfLaTeX (dos pasadas)
% =====================================================================
\documentclass[aspectratio=169,11pt,xcolor={table}]{beamer}

% ---------------------------------------------------------------------
%  Codificación, idioma y tipografía
% ---------------------------------------------------------------------
\usepackage[utf8]{inputenc}
\usepackage[T1]{fontenc}
\usepackage[spanish,es-nodecimaldot,es-noshorthands]{babel}
\usepackage{avant}                       % Sans geométrica (emula Avenir Next)
\renewcommand{\familydefault}{\sfdefault}
\renewcommand{\ttdefault}{pcr}           % Monoespaciada para código
\usepackage{textcomp}

% ---------------------------------------------------------------------
%  Paquetes
% ---------------------------------------------------------------------
\usepackage{amsmath,amssymb,mathtools}
\usepackage{siunitx}
\sisetup{output-decimal-marker={.}, per-mode=symbol}
\usepackage{booktabs,tabularx,multirow,array}
\usepackage{adjustbox}
\usepackage{tikz}
\usetikzlibrary{arrows.meta,positioning,calc,shapes.geometric,shapes.misc,
  decorations.pathreplacing,fit,backgrounds,babel}
\usepackage[siunitx,RPvoltages]{circuitikz}
\usepackage{pgfplots}
\pgfplotsset{compat=1.18}
\usepackage[most]{tcolorbox}
\usepackage{listings}

% ---------------------------------------------------------------------
%  Paleta institucional IUB (Manual de Identidad Visual 2025)
% ---------------------------------------------------------------------
\definecolor{iubwhite}{HTML}{FFFFFF}
\definecolor{iubnavy}{HTML}{121023}
\definecolor{iubyellow}{HTML}{FFDF2D}
\definecolor{iubblue}{HTML}{00ADE7}
\definecolor{iuborange}{HTML}{E39037}
\definecolor{iubgreen}{HTML}{3B9C6D}

% ---------------------------------------------------------------------
%  Configuración de Beamer
% ---------------------------------------------------------------------
\setbeamertemplate{navigation symbols}{}
\setbeamersize{text margin left=0.6cm, text margin right=0.6cm}
\setbeamercolor{normal text}{fg=iubnavy, bg=iubwhite}
\setbeamercolor{background canvas}{bg=iubwhite}
\setbeamercolor{structure}{fg=iubnavy}
\setbeamercolor{alerted text}{fg=iuborange}
\setbeamertemplate{itemize item}{\textcolor{iubblue}{\small$\blacktriangleright$}}
\setbeamertemplate{itemize subitem}{\textcolor{iuborange}{\scriptsize$\bullet$}}
\setbeamertemplate{enumerate item}{\textcolor{iubblue}{\bfseries\insertenumlabel.}}
\setbeamertemplate{frametitle}{}         % el título vive en el headline

% Parches: el headline se pre-mide antes del primer frame
\makeatletter
\providecommand{\insertframetitle}{}
\providecommand{\insertframesubtitle}{}
\makeatother

% ---------- Encabezado ----------
\setbeamertemplate{headline}{%
  \begin{tikzpicture}
    \useasboundingbox (0,0) rectangle (\paperwidth,1.05cm);
    \fill[iubnavy] (0,0) rectangle (\paperwidth,1.05cm);
    \fill[iubyellow] (0,0) rectangle (\paperwidth,0.05cm);
    \node[anchor=west, text=iubyellow, font=\large\bfseries]
      at (0.6cm,0.55cm) {\strut\insertframetitle};
    \node[anchor=east, text=iubyellow, font=\footnotesize\bfseries]
      at ([xshift=-0.6cm]\paperwidth,0.55cm) {IUB};
  \end{tikzpicture}}

% ---------- Pie de página ----------
\setbeamertemplate{footline}{%
  \begin{tikzpicture}
    \useasboundingbox (0,0) rectangle (\paperwidth,0.55cm);
    \fill[iubnavy] (0,0) rectangle (\paperwidth,0.55cm);
    \node[anchor=west, text=iubyellow, font=\tiny]
      at (0.6cm,0.275cm) {IEL04 \textperiodcentered{} Circuitos Eléctricos I};
    \node[anchor=center, text=iubyellow, font=\tiny]
      at (0.66\paperwidth,0.275cm) {Ing. Sergio Martinez Castilla};
    \node[anchor=east, text=iubyellow, font=\tiny\bfseries]
      at ([xshift=-0.6cm]\paperwidth,0.275cm)
      {Semana 8 \quad \insertframenumber\,/\,\inserttotalframenumber};
  \end{tikzpicture}}

% ---------------------------------------------------------------------
%  Cajas tcolorbox (texto blanco sobre la paleta complementaria)
%  \begin{caja}[opciones]{color}{título} ... \end{caja}
%  \begin{cajaplana}[opciones]{color} ... \end{cajaplana}
%  \begin{cardB|cardO|cardG|cardN}[opciones]{título} ... \end{cardX}
% ---------------------------------------------------------------------
\tcbset{
  iubbase/.style={enhanced, arc=1.5mm, boxrule=0pt, coltext=white, coltitle=white,
    fonttitle=\bfseries\footnotesize, fontupper=\footnotesize,
    left=1.8mm, right=1.8mm, top=1mm, bottom=1.2mm,
    toptitle=0.6mm, bottomtitle=0.6mm, halign=left,
    before upper={\hyphenpenalty=5000\setbeamercolor{itemize item}{fg=white}%
                  \setbeamercolor{itemize/enumerate body}{fg=white}%
                  \setbeamercolor{itemize/enumerate subbody}{fg=white}%
                  \setbeamercolor{itemize subitem}{fg=white}%
                  \setbeamercolor{enumerate item}{fg=white}%
                  \setbeamertemplate{itemize item}{\textcolor{white}{\small$\blacktriangleright$}}}}
}
\newtcolorbox{caja}[3][]{iubbase, colback=#2, colframe=#2,
  colbacktitle=#2!78!iubnavy, title={#3}, #1}
\newtcolorbox{cajaplana}[2][]{iubbase, colback=#2, colframe=#2, #1}
\newtcolorbox{cardB}[2][]{iubbase, colback=iubblue,   colframe=iubblue,
  colbacktitle=iubblue!78!iubnavy,   title={#2}, #1}
\newtcolorbox{cardO}[2][]{iubbase, colback=iuborange, colframe=iuborange,
  colbacktitle=iuborange!78!iubnavy, title={#2}, #1}
\newtcolorbox{cardG}[2][]{iubbase, colback=iubgreen,  colframe=iubgreen,
  colbacktitle=iubgreen!78!iubnavy,  title={#2}, #1}
\newtcolorbox{cardN}[2][]{iubbase, colback=iubnavy,   colframe=iubnavy,
  colbacktitle=iubnavy, coltitle=iubyellow, title={#2}, #1}

% Celda de encabezado de tabla (texto amarillo sobre \rowcolor{iubnavy}) y código en línea
\newcommand{\hd}[1]{\textcolor{iubyellow}{\bfseries #1}}
\newcommand{\thc}[1]{\textcolor{iubyellow}{\textbf{#1}}}
\newcommand{\cod}[1]{\texttt{\textbf{#1}}}

% ---------------------------------------------------------------------
%  Código sobre fondo oscuro:
%  \begin{codebox}[listing options app={language=Python,firstnumber=1}]{título}
% ---------------------------------------------------------------------
\lstdefinestyle{iubcode}{
  basicstyle=\ttfamily\fontsize{7}{8.4}\selectfont\color{white},
  keywordstyle=\color{iubblue}\bfseries,
  commentstyle=\color{iubgreen!55!white}\itshape,
  stringstyle=\color{iuborange!80!white},
  numbers=left, numberstyle=\tiny\color{iubyellow}, numbersep=6pt,
  showstringspaces=false, breaklines=true, tabsize=4,
  columns=fullflexible, keepspaces=true, upquote=true}
\newtcblisting{codebox}[2][]{enhanced, listing only, colback=iubnavy,
  colframe=iubnavy, colbacktitle=iubnavy!85!white, coltitle=iubyellow,
  fonttitle=\bfseries\scriptsize, title={#2}, arc=1.5mm, boxrule=0pt,
  left=6mm, right=1mm, top=0.5mm, bottom=0.5mm, toptitle=0.5mm, bottomtitle=0.5mm,
  listing options={style=iubcode}, #1}

% ---------------------------------------------------------------------
%  Estilos TikZ
% ---------------------------------------------------------------------
\tikzset{
  bloque/.style={rectangle, rounded corners=2mm, fill=#1, text=white,
    font=\scriptsize\bfseries, align=center, minimum height=1.1cm,
    text width=1.8cm, inner sep=1.2mm},
  flecha/.style={-{Stealth[length=2.2mm]}, line width=1pt, draw=iubnavy},
  arr/.style={flecha},
  term/.style={rectangle, draw=#1, line width=1.2pt, fill=white,
    text=iubnavy, font=\tiny\bfseries, minimum width=1.05cm,
    minimum height=0.5cm, inner sep=1pt},
  nodo/.style={rectangle, rounded corners=1mm, fill=#1, text=white,
    font=\scriptsize\bfseries, minimum size=0.75cm, align=center},
  cable/.style={line width=1.4pt, draw=#1},
  box/.style={rectangle, rounded corners=1.5mm, draw=none, text=white,
    font=\scriptsize\bfseries, align=center, minimum height=0.8cm, inner sep=3pt},
  bB/.style={box, fill=iubblue}, bO/.style={box, fill=iuborange},
  bG/.style={box, fill=iubgreen}, bN/.style={box, fill=iubnavy},
  dec/.style={diamond, aspect=1.9, fill=iuborange, text=white,
    font=\scriptsize\bfseries, align=center, inner sep=1pt},
  tag/.style={fill=#1, text=white, font=\scriptsize\bfseries, rounded corners=1mm,
    minimum width=0.7cm, anchor=north west},
}

% ---------------------------------------------------------------------
%  Diapositiva separadora de bloque
%  #1 número  #2 título  #3 duración  #4 color
% ---------------------------------------------------------------------
\newcommand{\bloque}[4]{%
\section{#2}
\begin{frame}[plain]
\begin{tikzpicture}[remember picture, overlay]
  \fill[#4] (current page.north west) rectangle
        ([xshift=5.2cm]current page.south west);
  \node[anchor=north west, text=white, font=\bfseries\large]
    at ([xshift=0.8cm,yshift=-2.2cm]current page.north west) {BLOQUE};
  \node[anchor=north west, text=white, font=\bfseries\fontsize{72}{72}\selectfont]
    at ([xshift=0.65cm,yshift=-2.9cm]current page.north west) {#1};
  \node[anchor=north west, text=iubnavy, font=\bfseries\LARGE,
        text width=9.4cm, align=left]
    at ([xshift=6.2cm,yshift=-2.8cm]current page.north west)
    {\hyphenpenalty=10000\exhyphenpenalty=10000 #2};
  \node[anchor=north west, fill=iubnavy, text=iubyellow, rounded corners=1.5mm,
        font=\small\bfseries, inner sep=2mm]
    at ([xshift=6.2cm,yshift=-6.0cm]current page.north west) {Duración: #3};
  \fill[iubnavy] ([xshift=5.2cm]current page.south west) rectangle
        ([yshift=0.35cm]current page.south east);
  \fill[iubyellow] ([xshift=5.2cm,yshift=0.35cm]current page.south west)
        rectangle ([yshift=0.42cm]current page.south east);
  \node[anchor=north east, text=iubnavy, font=\small\bfseries]
    at ([xshift=-0.6cm,yshift=-0.5cm]current page.north east) {IUB · IEL04};
\end{tikzpicture}
\end{frame}}

% =====================================================================
\begin{document}
% =====================================================================

% ---------------------------------------------------------------------
%  PORTADA
% ---------------------------------------------------------------------
\begin{frame}[plain]
\begin{tikzpicture}[remember picture, overlay]
  % Panel institucional
  \fill[iubnavy] (current page.north west) rectangle
        ([xshift=5.6cm]current page.south west);
  \node[anchor=north west, text=iubyellow, font=\bfseries\fontsize{54}{54}\selectfont]
    at ([xshift=0.7cm,yshift=-1.0cm]current page.north west) {IUB};
  \node[anchor=north west, text=white, font=\footnotesize, text width=4.3cm, align=left]
    at ([xshift=0.75cm,yshift=-3.0cm]current page.north west)
    {Institución Universitaria de Barranquilla};
  \draw[iubyellow, line width=1.2pt]
    ([xshift=0.75cm,yshift=-4.25cm]current page.north west) --
    ([xshift=2.6cm,yshift=-4.25cm]current page.north west);
  \node[anchor=north west, text=iubyellow, font=\scriptsize\bfseries,
        text width=4.3cm, align=left]
    at ([xshift=0.75cm,yshift=-4.5cm]current page.north west)
    {Tecnología en Gestión de Sistemas Eléctricos};
  \node[anchor=south west, text=white, font=\scriptsize, text width=4.3cm, align=left]
    at ([xshift=0.75cm,yshift=0.9cm]current page.south west)
    {Ing. Sergio Martinez Castilla\\[1pt]\textcolor{iubyellow}{\tiny smartinezc@unibarranquilla.edu.co}};
  % Bloque de título
  \node[anchor=north west, fill=iubblue, text=white, rounded corners=1.5mm,
        font=\scriptsize\bfseries, inner sep=2mm]
    at ([xshift=6.4cm,yshift=-1.0cm]current page.north west)
    {IEL04 \textperiodcentered{} Semana 8 de 14 \textperiodcentered{} Corte evaluativo 2};
  \node[anchor=north west, text=iubnavy, font=\small, text width=9.0cm, align=left] (a)
    at ([xshift=6.4cm,yshift=-1.9cm]current page.north west)
    {Circuitos Eléctricos I};
  \node[anchor=north west, text=iubnavy, font=\bfseries\fontsize{15}{18}\selectfont,
        text width=9.0cm, align=left] (t)
    at ([yshift=-0.35cm]a.south west)
    {\hyphenpenalty=10000 Circuitos RL en serie y en paralelo: reactancia inductiva, impedancia y leyes de Kirchhoff con la bobina real};
  % Franjas de la paleta complementaria
  \fill[iubblue]   ([xshift=6.4cm,yshift=1.1cm]current page.south west)
    rectangle ++(1.6cm,0.18cm);
  \fill[iuborange] ([xshift=8.1cm,yshift=1.1cm]current page.south west)
    rectangle ++(1.6cm,0.18cm);
  \fill[iubgreen]  ([xshift=9.8cm,yshift=1.1cm]current page.south west)
    rectangle ++(1.6cm,0.18cm);
  \node[anchor=north west, text=iubnavy, font=\scriptsize]
    at ([xshift=6.4cm,yshift=0.95cm]current page.south west)
    {Sesión teórico-práctica \textperiodcentered{} 240 min};
\end{tikzpicture}
\end{frame}

% ---------------------------------------------------------------------
\begin{frame}{Punto de partida de la sesión}
\begin{columns}[T,onlytextwidth]
  \column{0.34\textwidth}
\begin{caja}[height=5.9cm]{iubblue}{Continuidad}
\scriptsize \begin{itemize}\setlength{\itemsep}{2pt}
  \item Semana 7: Cálculo y medición de parámetros en circuitos RC serie
  \item \textbf{Semana 8: Circuitos RL en serie y en paralelo: tipos e identificación de resistores e inductores, análisis con las leyes de voltajes y corrientes de Kirchhoff}
  \item Semana 9: Cálculo y medición de parámetros en circuitos RL serie
\end{itemize}
\end{caja}
  \column{0.34\textwidth}
\begin{caja}[height=5.9cm]{iuborange}{Prerrequisitos}
\scriptsize \begin{itemize}\setlength{\itemsep}{2pt}
  \item Explicar la inductancia, la ley de Faraday y la resistencia
  \item Aplicar el método general de análisis fasorial de circuitos RC
  \item Calcular impedancias en serie y admitancias en paralelo en forma
  \item Identificar resistores por su código de colores e inductores
\end{itemize}
\end{caja}
  \column{0.29\textwidth}
\begin{caja}[height=5.9cm]{iubgreen}{Competencia}
\scriptsize \textbf{UC2.} Coordinar actividades asociadas a sistemas eléctricos utilizando fuentes convencionales y no convencionales de energía.
\end{caja}
\end{columns}
\end{frame}

% ---------------------------------------------------------------------
\begin{frame}{Resultados de aprendizaje}
\begin{columns}[c,onlytextwidth]
  \column{0.45\textwidth}
\begin{caja}{iubnavy}{IEL04-RA3}
\footnotesize Calcular un circuito resistivo, inductivo y RL, sea en serie o paralelo.
\end{caja}
\vspace{1mm}
\begin{cajaplana}{iubblue}
\scriptsize \textbf{CE1.} Identificar los tipos de resistores e inductores.\\[2pt]
\textbf{CE2.} Realizar mediciones en un circuito resistivo e inductivo.\\[2pt]
\textbf{CE3.} Realizar mediciones en un circuito RL.
\end{cajaplana}
  \column{0.52\textwidth}
\centering
\begin{adjustbox}{max width=\linewidth, max totalheight=6.1cm}
\begin{tikzpicture}
  \node[box, fill=iubblue, text width=4.9cm, align=left, font=\tiny, anchor=west] (r0) at (1.25,0.00) {Identificar el resistor y el inductor de un circuito RL y medir la resistencia de devanado que debe incluirse en el cálculo.};
  \node[tag=iubnavy, font=\tiny\bfseries, minimum width=1.15cm, anchor=east] at (1.1,0.00) {Aplicar};
  \node[box, fill=iuborange, text width=4.9cm, align=left, font=\tiny, anchor=west] (r1) at (1.25,-1.45) {Calcular la constante de tiempo y la reactancia inductiva de un circuito RL y explicar su dependencia con la frecuencia.};
  \node[tag=iubnavy, font=\tiny\bfseries, minimum width=1.15cm, anchor=east] at (1.1,-1.45) {Comprender};
  \node[box, fill=iubgreen, text width=4.9cm, align=left, font=\tiny, anchor=west] (r2) at (1.25,-2.90) {Analizar un circuito RL en serie y uno en paralelo aplicando las leyes de voltajes y de corrientes de Kirchhoff con fasores.};
  \node[tag=iubnavy, font=\tiny\bfseries, minimum width=1.15cm, anchor=east] at (1.1,-2.90) {Analizar};
  \node[box, fill=iubblue, text width=4.9cm, align=left, font=\tiny, anchor=west] (r3) at (1.25,-4.35) {Medir tensiones y corrientes en circuitos RL en serie y en paralelo y evaluar el efecto de la resistencia de devanado en los resultados.};
  \node[tag=iubnavy, font=\tiny\bfseries, minimum width=1.15cm, anchor=east] at (1.1,-4.35) {Evaluar};
  \draw[flecha, draw=iubnavy!40] (r0.south) -- (r1.north);
  \draw[flecha, draw=iubnavy!40] (r1.south) -- (r2.north);
  \draw[flecha, draw=iubnavy!40] (r2.south) -- (r3.north);
\end{tikzpicture}
\end{adjustbox}
\end{columns}
\end{frame}

% ---------------------------------------------------------------------
\begin{frame}{Ruta de la sesión \textperiodcentered{} 240 minutos}
\centering
\begin{tikzpicture}[x=1cm,y=1cm]
  \node[circle, fill=iubblue, text=white, font=\scriptsize\bfseries, minimum size=0.6cm, inner sep=0] at (0,0.00) {1};
  \node[anchor=west, font=\scriptsize, text width=7.6cm] at (0.45,0.00) {Qué cambia al sustituir el condensador por la bobina: reactancia creciente y tensión adelantada};
  \fill[iubblue] (8.3,-0.20) rectangle ++(2.20,0.4);
  \node[anchor=west, font=\scriptsize\bfseries] at (10.60,0.00) {20 min};
  \node[circle, fill=iuborange, text=white, font=\scriptsize\bfseries, minimum size=0.6cm, inner sep=0] at (0,-0.76) {2};
  \node[anchor=west, font=\scriptsize, text width=7.6cm] at (0.45,-0.76) {Tipos de resistores e inductores, resistencia de devanado y elección de componentes};
  \fill[iuborange] (8.3,-0.96) rectangle ++(2.75,0.4);
  \node[anchor=west, font=\scriptsize\bfseries] at (11.15,-0.76) {25 min};
  \node[circle, fill=iubgreen, text=white, font=\scriptsize\bfseries, minimum size=0.6cm, inner sep=0] at (0,-1.51) {3};
  \node[anchor=west, font=\scriptsize, text width=7.6cm] at (0.45,-1.51) {Crecimiento y caída exponencial de la corriente y \ensuremath{\tau} = L/R};
  \fill[iubgreen] (8.3,-1.71) rectangle ++(3.30,0.4);
  \node[anchor=west, font=\scriptsize\bfseries] at (11.70,-1.51) {30 min};
  \node[circle, fill=iubblue, text=white, font=\scriptsize\bfseries, minimum size=0.6cm, inner sep=0] at (0,-2.27) {4};
  \node[anchor=west, font=\scriptsize, text width=7.6cm] at (0.45,-2.27) {XL = 2\ensuremath{\pi}fL, tensión adelantada 90\textdegree{} y dependencia con la frecuencia};
  \fill[iubblue] (8.3,-2.47) rectangle ++(3.30,0.4);
  \node[anchor=west, font=\scriptsize\bfseries] at (11.70,-2.27) {30 min};
  \node[circle, fill=iuborange, text=white, font=\scriptsize\bfseries, minimum size=0.6cm, inner sep=0] at (0,-3.03) {5};
  \node[anchor=west, font=\scriptsize, text width=7.6cm] at (0.45,-3.03) {Z = R + jXL, ángulo de fase y suma fasorial VS = VR + VL};
  \fill[iuborange] (8.3,-3.23) rectangle ++(4.40,0.4);
  \node[anchor=west, font=\scriptsize\bfseries] at (12.80,-3.03) {40 min};
  \node[circle, fill=iubgreen, text=white, font=\scriptsize\bfseries, minimum size=0.6cm, inner sep=0] at (0,-3.79) {6};
  \node[anchor=west, font=\scriptsize, text width=7.6cm] at (0.45,-3.79) {Y = G \ensuremath{-} jBL, corrientes de rama y suma fasorial IT = IR + IL};
  \fill[iubgreen] (8.3,-3.99) rectangle ++(4.40,0.4);
  \node[anchor=west, font=\scriptsize\bfseries] at (12.80,-3.79) {40 min};
  \node[circle, fill=iubblue, text=white, font=\scriptsize\bfseries, minimum size=0.6cm, inner sep=0] at (0,-4.54) {7};
  \node[anchor=west, font=\scriptsize, text width=7.6cm] at (0.45,-4.54) {La bobina de 10.4 mH con un resistor de 680 \ensuremath{\Omega} a 10 kHz: efecto de la resistencia de devanado};
  \fill[iubblue] (8.3,-4.74) rectangle ++(4.40,0.4);
  \node[anchor=west, font=\scriptsize\bfseries] at (12.80,-4.54) {40 min};
  \node[circle, fill=iuborange, text=white, font=\scriptsize\bfseries, minimum size=0.6cm, inner sep=0] at (0,-5.30) {8};
  \node[anchor=west, font=\scriptsize, text width=7.6cm] at (0.45,-5.30) {Síntesis, cierre actitudinal y bibliografía};
  \fill[iuborange] (8.3,-5.50) rectangle ++(1.65,0.4);
  \node[anchor=west, font=\scriptsize\bfseries] at (10.05,-5.30) {15 min};
  \draw[iubnavy, line width=0.6pt] (8.3,0.45) -- (8.3,-5.70);
\end{tikzpicture}
\end{frame}

% =====================================================================
\bloque{01}{Qué cambia al sustituir el condensador por la bobina: reactancia creciente y tensión adelantada}{20 min}{iubblue}
% =====================================================================

% ---------------------------------------------------------------------
\begin{frame}{De los circuitos RC a los circuitos RL}
\begin{columns}[c,onlytextwidth]
  \column{0.57\textwidth}
\small
\begin{itemize}\setlength{\itemsep}{4pt}
  \item Si en cualquiera de los circuitos RC de las semanas anteriores se sustituye el condensador por una bobina, se obtiene un circuito RL.
  \item En la electrónica, los circuitos RL se usan como filtros, como limitadores de la corriente de arranque y como elementos de temporización.
  \item Los tres primeros son consecuencia directa de la física de la bobina: la reactancia inductiva, su crecimiento con la frecuencia y el adelanto de la tensión.
\end{itemize}
  \column{0.4\textwidth}
\begin{cardO}{Nota pedagógica}
\footnotesize El RA3 dispone solo de dos semanas antes del examen de la semana 10.
\end{cardO}
\end{columns}
\end{frame}

% ---------------------------------------------------------------------
\begin{frame}{{\normalsize Qué cambia al pasar de los circuitos RC a los circuitos RL}}
\centering
\begin{adjustbox}{max width=\linewidth, max totalheight=6.1cm}
\begin{tikzpicture}
  \node[bG, align=center, font=\scriptsize, text width=3.93cm] (n0) at (2.09,4.04) {Método general (semana 7)};
  \node[bB, align=center, font=\scriptsize, text width=2.84cm] (n1) at (7.34,4.04) {Sustituir C por L};
  \node[bB, align=center, font=\scriptsize, text width=2.84cm] (n2) at (12.05,4.04) {XL = 2\ensuremath{\pi}fL crece con f};
  \node[bB, align=center, font=\scriptsize, text width=2.84cm] (n3) at (12.05,1.01) {VL adelanta 90° a IL};
  \node[bO, align=center, font=\scriptsize, text width=3.26cm] (n4) at (7.34,1.01) {Incluir RW de la bobina};
  \node[bG, align=center, font=\scriptsize, text width=3.93cm] (n5) at (2.09,1.01) {RL serie (LVK) y paralelo (LCK)};
  \draw[flecha] (n0) -- (n1);
  \draw[flecha] (n1) -- (n2);
  \draw[flecha] (n2) -- (n3);
  \draw[flecha] (n3) -- (n4);
  \draw[flecha] (n4) -- (n5);
\end{tikzpicture}
\end{adjustbox}

\vspace{1mm}{\scriptsize Qué cambia al pasar de los circuitos RC a los circuitos RL\par}
\end{frame}

% =====================================================================
\bloque{02}{Tipos de resistores e inductores, resistencia de devanado y elección de componentes}{25 min}{iuborange}
% =====================================================================

% ---------------------------------------------------------------------
\begin{frame}{Resistores e inductores del circuito RL}
\begin{columns}[c,onlytextwidth]
  \column{0.57\textwidth}
\small
\begin{itemize}\setlength{\itemsep}{4pt}
  \item Para decidir si la resistencia de devanado importa en un circuito concreto, basta compararla con la reactancia inductiva a la frecuencia de trabajo.
\end{itemize}
  \column{0.4\textwidth}
\begin{cardO}{Error frecuente}
\footnotesize Medir la resistencia de devanado con el ohmímetro y restarla de la resistencia del circuito, o ignorarla por completo.
\end{cardO}
\end{columns}
\end{frame}

% ---------------------------------------------------------------------
\begin{frame}{{\normalsize Componentes del kit disponibles para los circuitos RL}}
\centering\tiny
\renewcommand{\arraystretch}{1.35}
\rowcolors{2}{iubnavy!6}{white}
\begin{adjustbox}{max width=\linewidth, max totalheight=5.6cm}
\begin{tabularx}{\textwidth}{>{\raggedright\arraybackslash}X>{\raggedright\arraybackslash}X>{\raggedright\arraybackslash}X>{\raggedright\arraybackslash}X>{\raggedright\arraybackslash}X>{\raggedright\arraybackslash}X>{\raggedright\arraybackslash}X}
  \rowcolor{iubnavy}
  \hd{Ref.} & \hd{Tipo} & \hd{Marcado} & \hd{Valor nominal} & \hd{Valor medido} & \hd{RW medida} & \hd{Q a 10 kHz}\\
  R & resistor de película & azul-gris-marrón-oro & 680 \ensuremath{\Omega} ±5 \% & 676 \ensuremath{\Omega} & — & —\\
  L1 & inductor radial & 102J & 1 mH ±5 \% & 1.03 mH & 3.2 \ensuremath{\Omega} & 20\\
  L2 & inductor moldeado & 472K & 4.7 mH ±10 \% & 4.62 mH & 9.8 \ensuremath{\Omega} & 30\\
  L3 & bobina con núcleo de ferrita & 103K & 10 mH ±10 \% & 10.4 mH & 24.6 \ensuremath{\Omega} & 27\\
\end{tabularx}
\end{adjustbox}
\vspace{2mm}
{\scriptsize Componentes del kit disponibles para los circuitos RL y su identificación\par}
\end{frame}

% =====================================================================
\bloque{03}{Crecimiento y caída exponencial de la corriente y \ensuremath{\tau} = L/R}{30 min}{iubgreen}
% =====================================================================

% ---------------------------------------------------------------------
\begin{frame}{Respuesta del circuito RL en serie con cd}
\begin{columns}[c,onlytextwidth]
  \column{0.57\textwidth}
\small
\begin{itemize}\setlength{\itemsep}{4pt}
  \item Al plantear la ley de voltajes de Kirchhoff, $V_S = iR + L\,di/dt$, se obtiene una ecuación diferencial de primer orden de la misma forma que la del RC.
  \item Con la bobina L3 del kit (10.4 mH) y el resistor de 676 \ensuremath{\Omega}, la resistencia total del lazo es $676 + 24.6 = 700.6\ \Omega$, y $\tau = 10.4 \times 10^{-3}/700.6 \approx 14.8\ \mu\text{s}$.
\end{itemize}
  \column{0.4\textwidth}
\begin{caja}{iubblue}{Ecuación clave}
\footnotesize \centering\small \begin{adjustbox}{max width=\linewidth}$\displaystyle \begin{gathered}i(t) = \frac{V_S}{R}\left(1 - e^{-t/\tau}\right) \\ v_L(t) = V_S\,e^{-t/\tau} \\ \tau = \frac{L}{R}\end{gathered}$\end{adjustbox}
\end{caja}
\end{columns}
\end{frame}

% ---------------------------------------------------------------------
\begin{frame}{{\normalsize Crecimiento y caída de la corriente en un RL en serie con VS}}
\begin{columns}[c,onlytextwidth]
  \column{0.62\textwidth}
\centering
\begin{tikzpicture}
  \begin{axis}[width=8.4cm, height=5.7cm, xmin=0, xmax=75, ymin=0, ymax=7.5684,
      xlabel={Tiempo t (µs)}, ylabel={Corriente i (mA)},
      label style={font=\scriptsize, text=iubnavy}, tick label style={font=\tiny, text=iubnavy},
      axis line style={iubnavy}, grid=major, grid style={iubnavy!12},
      legend style={font=\tiny, fill=white}, legend pos=north east]
    \addplot[line width=1.4pt, iubblue] coordinates {(0,0) (1.0033,0.468) (2.0067,0.90533) (3.01,1.314) (4.0134,1.6959) (5.0167,2.0527) (6.0201,2.3862) (7.0234,2.6978) (8.0268,2.9889) (9.0301,3.261) (10.033,3.5153) (11.037,3.7529) (12.04,3.9749) (13.043,4.1823) (14.047,4.3762) (15.05,4.5574) (16.054,4.7266) (17.057,4.8848) (18.06,5.0327) (19.064,5.1708) (20.067,5.2999) (21.07,5.4205) (22.074,5.5332) (23.077,5.6385) (24.08,5.7369) (25.084,5.8289) (26.087,5.9148) (27.09,5.9951) (28.094,6.0702) (29.097,6.1403) (30.1,6.2058) (31.104,6.2671) (32.107,6.3243) (33.11,6.3777) (34.114,6.4277) (35.117,6.4744) (36.12,6.518) (37.124,6.5588) (38.127,6.5969) (39.13,6.6325) (40.134,6.6658) (41.137,6.6968) (42.14,6.7259) (43.144,6.753) (44.147,6.7784) (45.151,6.8021) (46.154,6.8242) (47.157,6.8449) (48.161,6.8643) (49.164,6.8824) (50.167,6.8992) (51.171,6.915) (52.174,6.9298) (53.177,6.9435) (54.181,6.9564) (55.184,6.9685) (56.187,6.9797) (57.191,6.9902) (58.194,7) (59.197,7.0092) (60.201,7.0178) (61.204,7.0258) (62.207,7.0333) (63.211,7.0403) (64.214,7.0468) (65.217,7.0529) (66.221,7.0586) (67.224,7.064) (68.227,7.0689) (69.231,7.0736) (70.234,7.078) (71.237,7.082) (72.241,7.0858) (73.244,7.0894) (74.247,7.0927) (75,7.095)};
    \addlegendentry{Crecimiento al conectar}
    \addplot[line width=1.4pt, iuborange] coordinates {(0,7.14) (1.0033,6.672) (2.0067,6.2347) (3.01,5.826) (4.0134,5.4441) (5.0167,5.0873) (6.0201,4.7538) (7.0234,4.4422) (8.0268,4.1511) (9.0301,3.879) (10.033,3.6247) (11.037,3.3871) (12.04,3.1651) (13.043,2.9577) (14.047,2.7638) (15.05,2.5826) (16.054,2.4134) (17.057,2.2552) (18.06,2.1073) (19.064,1.9692) (20.067,1.8401) (21.07,1.7195) (22.074,1.6068) (23.077,1.5015) (24.08,1.4031) (25.084,1.3111) (26.087,1.2252) (27.09,1.1449) (28.094,1.0698) (29.097,0.9997) (30.1,0.93417) (31.104,0.87294) (32.107,0.81572) (33.11,0.76225) (34.114,0.71229) (35.117,0.6656) (36.12,0.62198) (37.124,0.58121) (38.127,0.54311) (39.13,0.50751) (40.134,0.47425) (41.137,0.44316) (42.14,0.41411) (43.144,0.38697) (44.147,0.3616) (45.151,0.3379) (46.154,0.31575) (47.157,0.29506) (48.161,0.27572) (49.164,0.25765) (50.167,0.24076) (51.171,0.22498) (52.174,0.21023) (53.177,0.19645) (54.181,0.18357) (55.184,0.17154) (56.187,0.1603) (57.191,0.14979) (58.194,0.13997) (59.197,0.1308) (60.201,0.12222) (61.204,0.11421) (62.207,0.10673) (63.211,0.099731) (64.214,0.093194) (65.217,0.087085) (66.221,0.081377) (67.224,0.076043) (68.227,0.071059) (69.231,0.066401) (70.234,0.062049) (71.237,0.057982) (72.241,0.054181) (73.244,0.05063) (74.247,0.047311) (75,0.044966)};
    \addlegendentry{Caída al cortocircuitar la fuente}
    \addplot[line width=1pt, dashed, iubnavy!70] coordinates {(14.8,0) (14.8,7.5684)};
    \addlegendentry{t = \ensuremath{\tau}}
    \addplot[line width=1pt, dashed, iubnavy!70] coordinates {(0,4.5) (75,4.5)};
    \addlegendentry{63 \% de Imax}
  \end{axis}
\end{tikzpicture}
  \column{0.35\textwidth}
\begin{caja}{iubblue}{Lectura de la gráfica}
\footnotesize La gráfica es idéntica en forma a la de carga y descarga del condensador, con la corriente en lugar de la tensión.
\end{caja}
\end{columns}
\end{frame}

% =====================================================================
\bloque{04}{XL = 2\ensuremath{\pi}fL, tensión adelantada 90\textdegree{} y dependencia con la frecuencia}{30 min}{iubblue}
% =====================================================================

% ---------------------------------------------------------------------
\begin{frame}{{\normalsize Reactancia inductiva y respuesta sinusoidal de la bobina}}
\begin{columns}[c,onlytextwidth]
  \column{0.57\textwidth}
\small
\begin{itemize}\setlength{\itemsep}{4pt}
  \item Con una fuente sinusoidal, la corriente de la bobina cambia continuamente, y con ella la tensión inducida.
  \item La oposición de la bobina a la corriente sinusoidal se llama \textbf{reactancia inductiva}, se mide en ohmios $[1]$ y se simboliza $X_L$.
  \item Los extremos de la fórmula confirman lo ya visto.
\end{itemize}
  \column{0.4\textwidth}
\begin{caja}{iubblue}{Ecuación clave}
\footnotesize \centering\small \begin{adjustbox}{max width=\linewidth}$\displaystyle X_L = 2\pi f L$\end{adjustbox}
\end{caja}
\end{columns}
\end{frame}

% ---------------------------------------------------------------------
\begin{frame}{{\normalsize Reactancia inductiva de la bobina L3 (10.4 mH) en función}}
\begin{columns}[c,onlytextwidth]
  \column{0.62\textwidth}
\centering
\begin{tikzpicture}
  \begin{axis}[width=8.4cm, height=5.7cm, xmin=0, xmax=20000, ymin=0, ymax=1385.3,
      xlabel={Frecuencia f (Hz)}, ylabel={Reactancia XL (\ensuremath{\Omega})},
      label style={font=\scriptsize, text=iubnavy}, tick label style={font=\tiny, text=iubnavy},
      axis line style={iubnavy}, grid=major, grid style={iubnavy!12},
      legend style={font=\tiny, fill=white}, legend pos=north east]
    \addplot[line width=1.4pt, iubblue] coordinates {(0,0) (267.56,17.484) (535.12,34.967) (802.68,52.451) (1070.2,69.935) (1337.8,87.418) (1605.4,104.9) (1872.9,122.39) (2140.5,139.87) (2408,157.35) (2675.6,174.84) (2943.1,192.32) (3210.7,209.8) (3478.3,227.29) (3745.8,244.77) (4013.4,262.25) (4280.9,279.74) (4548.5,297.22) (4816.1,314.71) (5083.6,332.19) (5351.2,349.67) (5618.7,367.16) (5886.3,384.64) (6153.8,402.12) (6421.4,419.61) (6689,437.09) (6956.5,454.57) (7224.1,472.06) (7491.6,489.54) (7759.2,507.03) (8026.8,524.51) (8294.3,541.99) (8561.9,559.48) (8829.4,576.96) (9097,594.44) (9364.5,611.93) (9632.1,629.41) (9899.7,646.89) (10167,664.38) (10435,681.86) (10702,699.35) (10970,716.83) (11237,734.31) (11505,751.8) (11773,769.28) (12040,786.76) (12308,804.25) (12575,821.73) (12843,839.22) (13110,856.7) (13378,874.18) (13645,891.67) (13913,909.15) (14181,926.63) (14448,944.12) (14716,961.6) (14983,979.08) (15251,996.57) (15518,1014.1) (15786,1031.5) (16054,1049) (16321,1066.5) (16589,1084) (16856,1101.5) (17124,1119) (17391,1136.4) (17659,1153.9) (17926,1171.4) (18194,1188.9) (18462,1206.4) (18729,1223.9) (18997,1241.3) (19264,1258.8) (19532,1276.3) (19799,1293.8) (20000,1306.9)};
    \addlegendentry{XL}
    \addplot[line width=1pt, dashed, iubnavy!70] coordinates {(0,700.6) (20000,700.6)};
    \addlegendentry{R + RW = 700.6 \ensuremath{\Omega}}
    \addplot[line width=1pt, dashed, iubnavy!70] coordinates {(10722,0) (10722,1385.3)};
    \addlegendentry{XL = R (10.7 kHz)}
  \end{axis}
\end{tikzpicture}
  \column{0.35\textwidth}
\begin{caja}{iubblue}{Lectura de la gráfica}
\footnotesize La recta cruza la resistencia total en unos 10.7 kHz: por debajo de esa frecuencia domina el resistor, y por encima, la bobina.
\end{caja}
\end{columns}
\end{frame}

% =====================================================================
\bloque{05}{Z = R + jXL, ángulo de fase y suma fasorial VS = VR + VL}{40 min}{iuborange}
% =====================================================================

% ---------------------------------------------------------------------
\begin{frame}{{\normalsize Circuito RL en serie y ley de voltajes de Kirchhoff}}
\begin{columns}[c,onlytextwidth]
  \column{0.57\textwidth}
\small
\begin{itemize}\setlength{\itemsep}{4pt}
  \item En un circuito RL en serie, la corriente es común a ambos elementos y se toma como referencia.
  \item La ley de voltajes de Kirchhoff se aplica, como en el RC, sumando fasores.
\end{itemize}
  \column{0.4\textwidth}
\begin{caja}{iubblue}{Ecuación clave}
\footnotesize \centering\small \begin{adjustbox}{max width=\linewidth}$\displaystyle \mathbf{Z} = R + jX_L = \sqrt{R^{2} + X_L^{2}}\ \angle \tan^{-1}\!\left(\frac{X_L}{R}\right)$\end{adjustbox}
\end{caja}
\end{columns}
\end{frame}

% ---------------------------------------------------------------------
\begin{frame}{Ecuaciones: Circuito RL en serie y ley de voltajes}
\begin{columns}[c,onlytextwidth]
  \column{0.55\textwidth}
\begin{cajaplana}{iubblue}
\centering\small \begin{adjustbox}{max width=\linewidth}$\displaystyle \begin{gathered}\mathbf{V}_S = V_R + jV_L \\ V_S = \sqrt{V_R^{2} + V_L^{2}} \\ \theta = \tan^{-1}\!\left(\frac{V_L}{V_R}\right)\end{gathered}$\end{adjustbox}
\end{cajaplana}
  \column{0.42\textwidth}
\begin{cardO}{Error frecuente}
\footnotesize Usar el signo del circuito RC y escribir $\mathbf{Z} = R - jX_L$.
\end{cardO}
\end{columns}
\end{frame}

% ---------------------------------------------------------------------
\begin{frame}{{\normalsize Impedancia y ángulo de fase de un RL en serie ideal}}
\centering\scriptsize
\renewcommand{\arraystretch}{1.35}
\rowcolors{2}{iubnavy!6}{white}
\begin{adjustbox}{max width=\linewidth, max totalheight=5.6cm}
\begin{tabularx}{\textwidth}{>{\raggedright\arraybackslash}X>{\raggedright\arraybackslash}X>{\raggedright\arraybackslash}X>{\raggedright\arraybackslash}X>{\raggedright\arraybackslash}X}
  \rowcolor{iubnavy}
  \hd{f (Hz)} & \hd{XL (\ensuremath{\Omega})} & \hd{Z (\ensuremath{\Omega})} & \hd{\ensuremath{\theta} (°)} & \hd{Comportamiento}\\
  1000 & 62.8 & 683 & 5.3 & casi resistivo\\
  5000 & 314 & 749 & 24.8 & predomina R\\
  10 000 & 628 & 926 & 42.7 & equilibrado\\
  10 823 & 680 & 962 & 45.0 & XL = R\\
  20 000 & 1257 & 1429 & 61.6 & predomina L\\
\end{tabularx}
\end{adjustbox}
\vspace{2mm}
{\scriptsize Impedancia y ángulo de fase de un RL en serie ideal (R = 680 \ensuremath{\Omega}, L = 10 mH) en función de la frecuencia\par}
\end{frame}

% =====================================================================
\bloque{06}{Y = G \ensuremath{-} jBL, corrientes de rama y suma fasorial IT = IR + IL}{40 min}{iubgreen}
% =====================================================================

% ---------------------------------------------------------------------
\begin{frame}{{\normalsize Circuito RL en paralelo y ley de corrientes de Kirchhoff}}
\begin{columns}[c,onlytextwidth]
  \column{0.57\textwidth}
\small
\begin{itemize}\setlength{\itemsep}{4pt}
  \item En el circuito RL en paralelo, la tensión de la fuente es común al resistor y a la bobina, y se toma como referencia.
\end{itemize}
  \column{0.4\textwidth}
\begin{caja}{iubblue}{Ecuación clave}
\footnotesize \centering\small \begin{adjustbox}{max width=\linewidth}$\displaystyle \begin{gathered}\mathbf{Y}= G - jB_L \\ = \sqrt{G^{2} + B_L^{2}}\ \angle -\tan^{-1}\!\left(\frac{B_L}{G}\right) \\ B_L = \frac{1}{2\pi f L}\end{gathered}$\end{adjustbox}
\end{caja}
\end{columns}
\end{frame}

% ---------------------------------------------------------------------
\begin{frame}{{\normalsize Ecuaciones: Circuito RL en paralelo y ley de corrientes}}
\begin{columns}[c,onlytextwidth]
  \column{0.55\textwidth}
\begin{cajaplana}{iubblue}
\centering\small \begin{adjustbox}{max width=\linewidth}$\displaystyle \begin{gathered}I_T = \sqrt{I_R^{2} + I_L^{2}} \\ \theta = -\tan^{-1}\!\left(\frac{I_L}{I_R}\right)\end{gathered}$\end{adjustbox}
\end{cajaplana}
  \column{0.42\textwidth}
\begin{cardO}{Advertencia de seguridad}
\footnotesize Nunca se conecta un RL en paralelo directamente a una fuente de cd o de muy baja frecuencia sin una resistencia en serie que limite la corriente: la bobina se comporta como un cortocircuito y puede dañar la fuente, el generador o la propia bobina por sobrecalentamiento.
\end{cardO}
\end{columns}
\end{frame}

% ---------------------------------------------------------------------
\begin{frame}{Corrientes de un RL en paralelo ideal}
\centering\tiny
\renewcommand{\arraystretch}{1.35}
\rowcolors{2}{iubnavy!6}{white}
\begin{adjustbox}{max width=\linewidth, max totalheight=5.6cm}
\begin{tabularx}{\textwidth}{>{\raggedright\arraybackslash}X>{\raggedright\arraybackslash}X>{\raggedright\arraybackslash}X>{\raggedright\arraybackslash}X>{\raggedright\arraybackslash}X>{\raggedright\arraybackslash}X}
  \rowcolor{iubnavy}
  \hd{f (Hz)} & \hd{XL (\ensuremath{\Omega})} & \hd{IR (mA)} & \hd{IL (mA)} & \hd{IT (mA)} & \hd{\ensuremath{\theta} (°)}\\
  1000 & 62.8 & 7.35 & 79.6 & 79.9 & \ensuremath{-}84.7\\
  5000 & 314 & 7.35 & 15.9 & 17.5 & \ensuremath{-}65.2\\
  10 000 & 628 & 7.35 & 7.96 & 10.8 & \ensuremath{-}47.3\\
  20 000 & 1257 & 7.35 & 3.98 & 8.36 & \ensuremath{-}28.4\\
\end{tabularx}
\end{adjustbox}
\vspace{2mm}
{\scriptsize Corrientes de un RL en paralelo ideal (R = 680 \ensuremath{\Omega}, L = 10 mH, VS = 5 V eficaces) en función de la frecuencia\par}
\end{frame}

% =====================================================================
\bloque{07}{La bobina de 10.4 mH con un resistor de 680 \ensuremath{\Omega} a 10 kHz: efecto de la resistencia de devanado}{40 min}{iubblue}
% =====================================================================

% ---------------------------------------------------------------------
\begin{frame}{Medición de circuitos RL en serie y en paralelo}
\begin{columns}[c,onlytextwidth]
  \column{0.57\textwidth}
\small
\begin{itemize}\setlength{\itemsep}{4pt}
  \item Con $X_L = 2\pi(10\ 000)(0.0104) \approx 653.5\ \Omega$ y la resistencia total $676 + 24.6 = 700.6\ \Omega$, la impedancia es $\mathbf{Z} = 700.6 + j653.5 \approx 958 \angle 43.0^\circ\ \Omega$ y la corriente $I = 5/958 \approx 5.22\ \text{mA}$.
  \item Con los mismos componentes en paralelo y la misma tensión: $I_R = 5/676 \approx 7.40\ \text{mA}$ e $I_{bob} = 5/654 \approx 7.65\ \text{mA}$, retrasada 87.8° respecto a la tensión.
  \item La columna del cálculo ideal permite ver dónde pesa la resistencia de devanado.
\end{itemize}
  \column{0.4\textwidth}
\begin{caja}{iubblue}{Ecuación clave}
\footnotesize \centering\small \begin{adjustbox}{max width=\linewidth}$\displaystyle \begin{gathered}\mathbf{V}_S = V_R\angle 0^\circ + V_{bob}\angle \varphi_b \\ \varphi_b = \tan^{-1}\!\left(\frac{X_L}{R_W}\right)\end{gathered}$\end{adjustbox}
\end{caja}
\end{columns}
\end{frame}

% ---------------------------------------------------------------------
\begin{frame}{Circuitos RL en serie y en paralelo a 10 kHz (1/2)}
\centering\scriptsize
\renewcommand{\arraystretch}{1.35}
\rowcolors{2}{iubnavy!6}{white}
\begin{adjustbox}{max width=\linewidth, max totalheight=5.6cm}
\begin{tabularx}{\textwidth}{>{\raggedright\arraybackslash}X>{\raggedright\arraybackslash}X>{\raggedright\arraybackslash}X>{\raggedright\arraybackslash}X>{\raggedright\arraybackslash}X}
  \rowcolor{iubnavy}
  \hd{Magnitud} & \hd{Cálculo con RW} & \hd{Cálculo ideal (sin RW)} & \hd{Medido} & \hd{Diferencia con el cálculo con RW}\\
  Serie: I & 5.22 mA & 5.32 mA & 5.21 mA & \ensuremath{-}0.2 \%\\
  Serie: VR & 3.53 V & 3.60 V & 3.52 V & \ensuremath{-}0.3 \%\\
  Serie: Vbob & 3.41 V & 3.48 V & 3.42 V & +0.3 \%\\
  Serie: \ensuremath{\theta} & 43.0° & 44.0° & 43.2° & 0.2°\\
\end{tabularx}
\end{adjustbox}
\vspace{2mm}
{\scriptsize Circuitos RL en serie y en paralelo a 10 kHz con la bobina L3: cálculo con bobina real, cálculo ideal y medición (lecturas de ejemplo)\par}
\end{frame}

% ---------------------------------------------------------------------
\begin{frame}{Circuitos RL en serie y en paralelo a 10 kHz (2/2)}
\centering\scriptsize
\renewcommand{\arraystretch}{1.35}
\rowcolors{2}{iubnavy!6}{white}
\begin{adjustbox}{max width=\linewidth, max totalheight=5.6cm}
\begin{tabularx}{\textwidth}{>{\raggedright\arraybackslash}X>{\raggedright\arraybackslash}X>{\raggedright\arraybackslash}X>{\raggedright\arraybackslash}X>{\raggedright\arraybackslash}X}
  \rowcolor{iubnavy}
  \hd{Magnitud} & \hd{Cálculo con RW} & \hd{Cálculo ideal (sin RW)} & \hd{Medido} & \hd{Diferencia con el cálculo con RW}\\
  Paralelo: IR & 7.40 mA & 7.40 mA & 7.38 mA & \ensuremath{-}0.3 \%\\
  Paralelo: Ibob & 7.65 mA & 7.65 mA & 7.66 mA & +0.1 \%\\
  Paralelo: IT & 10.8 mA & 10.6 mA & 10.8 mA & 0.0 \%\\
  Paralelo: \ensuremath{\theta} & \ensuremath{-}44.8° & \ensuremath{-}46.0° & \ensuremath{-}44.9° & 0.1°\\
\end{tabularx}
\end{adjustbox}
\vspace{2mm}
{\scriptsize Circuitos RL en serie y en paralelo a 10 kHz con la bobina L3: cálculo con bobina real, cálculo ideal y medición (lecturas de ejemplo)\par}
\end{frame}

% =====================================================================
\bloque{08}{Síntesis, cierre actitudinal y bibliografía}{15 min}{iuborange}
% =====================================================================

% ---------------------------------------------------------------------
\begin{frame}{Síntesis de la sesión}
\begin{columns}[T,onlytextwidth]
  \column{0.32\textwidth}
\begin{cardB}{Resistores e inductores del circuito RL}
\scriptsize Tipos de resistores e inductores, resistencia de devanado y elección de componentes
\end{cardB}
  \column{0.32\textwidth}
\begin{cardO}{Respuesta del circuito RL en serie con cd}
\scriptsize Crecimiento y caída exponencial de la corriente y \ensuremath{\tau} = L/R
\end{cardO}
  \column{0.32\textwidth}
\begin{cardG}{Reactancia inductiva y respuesta}
\scriptsize XL = 2\ensuremath{\pi}fL, tensión adelantada 90° y dependencia con la frecuencia
\end{cardG}
\end{columns}
\vspace{2mm}
\begin{columns}[T,onlytextwidth]
  \column{0.32\textwidth}
\begin{cardB}{Circuito RL en serie y ley de voltajes}
\scriptsize Z = R + jXL, ángulo de fase y suma fasorial VS = VR + VL
\end{cardB}
  \column{0.32\textwidth}
\begin{cardO}{Circuito RL en paralelo y ley}
\scriptsize Y = G \ensuremath{-} jBL, corrientes de rama y suma fasorial IT = IR + IL
\end{cardO}
  \column{0.32\textwidth}
\begin{cardG}{Medición de circuitos RL en serie}
\scriptsize La bobina de 10.4 mH con un resistor de 680 \ensuremath{\Omega} a 10 kHz: efecto de la resistencia de devanado
\end{cardG}
\end{columns}
\end{frame}

% ---------------------------------------------------------------------
\begin{frame}{Reflexión para llevar}
\centering
\begin{tikzpicture}
  \node[fill=iubnavy, text=white, rounded corners=6pt, text width=11.5cm, inner sep=12pt, align=center, font=\normalsize] (c) at (0,0) {\itshape «Advertir al compañero antes de conectar una bobina en paralelo a baja frecuencia, o de abrir un circuito inductivo con corriente, es parte del trabajo en equipo responsable: en los circuitos RL, un descuido de conexión se paga con un generador dañado o una chispa en la mano.»};
  \node[fill=iubyellow, text=iubnavy, rounded corners=3pt, font=\scriptsize\bfseries, inner sep=4pt, anchor=south west] at ([yshift=-4pt]c.north west) {Componente actitudinal};
\end{tikzpicture}
\end{frame}

% ---------------------------------------------------------------------
\begin{frame}{Referencias bibliográficas}
\small
\begin{itemize}\setlength{\itemsep}{8pt}
  \item[{$[1]$}] T. L. Floyd, Principios de circuitos eléctricos, 8.\textordfeminine{} ed. México: Pearson Educación, 2007.
  \item[{$[2]$}] R. L. Boylestad, Introducción al análisis de circuitos, 10.\textordfeminine{} ed. México: Pearson Educación, 2004.
  \item[{$[3]$}] M. F. P. Deorsola y P. Morcelle del Valle, Circuitos eléctricos. Parte 1, 1.\textordfeminine{} ed. La Plata: Editorial de la Universidad de La Plata, 2017.
\end{itemize}
\end{frame}

% ---------------------------------------------------------------------
%  CIERRE
% ---------------------------------------------------------------------
\begin{frame}[plain]
\begin{tikzpicture}[remember picture, overlay]
  \fill[iubnavy] (current page.north west) rectangle (current page.south east);
  \node[anchor=north west, text=iubyellow, font=\bfseries\fontsize{40}{44}\selectfont]
    at ([xshift=1.2cm,yshift=-1.6cm]current page.north west) {\textexclamdown{}Gracias!};
  \node[anchor=north west, text=white, font=\normalsize, text width=14cm, align=left]
    at ([xshift=1.2cm,yshift=-3.4cm]current page.north west)
    {IEL04 \textperiodcentered{} Circuitos Eléctricos I};
  \node[anchor=north west, fill=iubblue, text=white, rounded corners=1.5mm,
        font=\small\bfseries, inner sep=2.2mm, text width=11.5cm]
    at ([xshift=1.2cm,yshift=-4.4cm]current page.north west)
    {Próxima sesión \textperiodcentered{} Semana 9: Cálculo y medición de parámetros en circuitos RL serie y paralelo según el tipo de conexionado};
  \node[anchor=south west, text=iubyellow, font=\small, align=left]
    at ([xshift=1.2cm,yshift=0.9cm]current page.south west)
    {Ing. Sergio Martinez Castilla \quad · \quad smartinezc@unibarranquilla.edu.co};
  \fill[iubblue]   ([xshift=1.2cm,yshift=0.55cm]current page.south west) rectangle ++(1.6cm,0.15cm);
  \fill[iuborange] ([xshift=2.9cm,yshift=0.55cm]current page.south west) rectangle ++(1.6cm,0.15cm);
  \fill[iubgreen]  ([xshift=4.6cm,yshift=0.55cm]current page.south west) rectangle ++(1.6cm,0.15cm);
\end{tikzpicture}
\end{frame}

\end{document}
